\begin{abcliste}
\abc In the wave model, a brighter light is a stronger wave: it shakes the electrons harder. It predicts that the $\result{\text{stopping voltage rises}}$ (faster electrons) and that the $\result{\text{saturation current rises}}$ (more energy, more electrons).

\abc The saturation current grows in proportion to the power,
\begin{align}
 I_\mathrm{s}' &= \IbF = \kf\times\Is = \Ib \approx \IbP{-6}{3}
\end{align}
as both models expect: that prediction $\result{\text{holds}}$. The stopping voltage stays the same: the fastest electrons are not faster in brighter light. The wave model's prediction for the stopping voltage $\result{\text{fails}}$.

\abc Each photon brings
\begin{align}
 E_\mathrm{ph} &= \EpVF = \EpV \approx \result{\EpVP{0}{3}}
\end{align}
After paying the work function, the fastest electrons keep
\begin{align}
 E_\mathrm{k,max} &= \EkF = \EpV - \WaO = \Ek \approx \result{\EkP{0}{3}}
\end{align}
which a counter-voltage of $\SI{0.59}{\V}$ just stops, as measured. This energy depends only on the photon, that is on the wavelength, not on how many photons there are.

Photons and electrons per second:
\begin{align}
 N_\mathrm{ph} &= \NaF = \Na \approx \result{\NanP{0}{2}\cdot 10^{15}\,\si{\per\second}} & N_e &= \NeF = \Ne \approx \result{\NenP{0}{2}\cdot 10^{13}\,\si{\per\second}} \\
 N_\mathrm{ph}' &= \Nb & N_e' &= \Neb
\end{align}
Only about $\SI{1.3}{\percent}$ of the photons release an electron ($N_e/N_\mathrm{ph} \approx \yaP{0}{2}$), but this share stays the same: $\kfO$ times as many photons, each with the same energy, release $\kfO$ times as many electrons, each with the same largest kinetic energy. Brighter light gives more electrons, not faster ones.
\end{abcliste}
